\section{Функцийн шаардлагууд}

\begin{enumerate}
    \item \textbf{FR-1}: Систем нь хэрэглэгчийг цахим шуудан, нууц үг ашиглан бүртгэх, нэвтрэх, нууц үг сэргээх боломжтой байх ёстой.
    \item \textbf{FR-2}: Систем нь хэрэглэгчийн эрхийн түвшингээр нь ялгаж, хэрэглэгч болон админы интерфэйсийг тусад нь харуулах ёстой.
    \item \textbf{FR-3}: Систем нь хэрэглэгчид шинэ хичээл үүсгэх, хичээлийн нэр, тайлбар, нууцлалын төлөв тохируулах боломж олгох ёстой.
    \item \textbf{FR-4}: Систем нь PDF болон зурган файл хичээл эсвэл тэмдэглэлд хавсаргах боломжтой байх ёстой.
    \item \textbf{FR-5}: Систем нь хэрэглэгчийн оруулсан номын материалыг автоматаар боловсруулж, хичээл $\rightarrow$ сэдэв $\rightarrow$ дэд сэдэв гэсэн бүтэцтэй хэсгүүдэд хуваах ёстой.
    \item \textbf{FR-6}: Систем нь сканнердсан зураг болон текст агуулаагүй PDF-ээс OCR ашиглан текст таних ёстой.
    \item \textbf{FR-7}: Систем танилт хийсэн текстийг эмбед хийн вектор өгөгдлийн сан (Vector Database)-д хадгалах
    \item \textbf{FR-8}: Систем нь материалын агуулга дээр үндэслэн AI-аар хураангуй, түлхүүр ойлголт, флаш карт, сорилын асуулт автоматаар үүсгэх ёстой.
    \item \textbf{FR-9}: Систем нь хэрэглэгчид AI-аар үүссэн контентыг засварлах, гараар нэмэлт тэмдэглэл оруулах, хүссэн хэсгээ устгах боломж олгох ёстой.
    \item \textbf{FR-10}: Систем нь хэрэглэгч чатботоор тухайн хичээлийн агуулгаар асуулт асуухад хэрэглэгчийн асуултыг эмбед хийн вектор өгөгдлийн сангаас хайлт хийн үр дүнд суурилан AI ашиглан хариулт өгнө.
    \item \textbf{FR-11}: Систем нь хэрэглэгчид тест өгөх, зөв бурууг нь харах, оноогоо хадгалах боломжтой байх.
    \item \textbf{FR-12}: Систем нь хэрэглэгчид хичээлээ Public эсвэл Private төлөвтэй болгох, Public үед хуваалцах холбоос үүсгэх боломж олгох ёстой.
    \item \textbf{FR-13}: Систем нь админд хэрэглэгч, хичээл, AI боловсруулалтын үр дүн, нийт статистик мэдээллийг харах, шаардлагатай тохиолдолд агуулгыг дахин боловсруулах боломж олгох ёстой.
\end{enumerate}

\section{Функцийн бус шаардлагууд}

\subsection{Гүйцэтгэл}

\begin{table}[H]
\centering
\begin{tabular}{|p{3cm}|p{11cm}|}
\hline
\textbf{ID} & \textbf{Шаардлага} \\ \hline
NFR-P1 & Систем нь энгийн хуудасны ачааллыг ихэнх тохиолдолд 3 секундээс бага хугацаанд гүйцэтгэх ёстой. \\ \hline
NFR-P2 & Файл оруулсны дараах анхны хариу болон боловсруулалт эхэлсэн төлөвийг хэрэглэгчид 5 секундээс бага хугацаанд харуулах ёстой. \\ \hline
NFR-P3 & Дунд хэмжээний материалд AI болон OCR боловсруулалтын гол үр дүнг 30 секунд орчим хугацаанд гаргах. \\ \hline
\end{tabular}
\caption{Гүйцэтгэлийн функцийн бус шаардлагууд}
\end{table}

\subsection{Аюулгүй байдал}

\begin{table}[H]
\centering
\begin{tabular}{|p{3cm}|p{11cm}|}
\hline
\textbf{ID} & \textbf{Шаардлага} \\ \hline
NFR-S1 & Систем нь HTTPS/TLS ашиглан клиент болон серверийн хооронд дамжиж буй өгөгдлийг шифрлэх ёстой. \\ \hline
NFR-S2 & Хэрэглэгчийн нууц үгийг аюулгүй хэш алгоритмаар хадгалах ёстой. \\ \hline
NFR-S3 & Нэвтрэлтийн төлвийг server-side session эсвэл түүнтэй адилтгах token механизмаар удирдах ёстой. \\ \hline
NFR-S4 & Админ болон энгийн хэрэглэгчийн хандалтыг role-based access control механизмаар ялгах ёстой. \\ \hline
\end{tabular}
\caption{Аюулгүй байдлын функцийн бус шаардлагууд}
\end{table}

\subsection{Өргөтгөх чадвар}

\begin{table}[H]
\centering
\begin{tabular}{|p{3cm}|p{11cm}|}
\hline
\textbf{ID} & \textbf{Шаардлага} \\ \hline
NFR-SC1 & Систем нь зэрэгцээ олон хэрэглэгчийн ердийн ачааллыг гүйцэтгэлийн үндсэн шаардлагыг зөрчихгүйгээр дэмжих ёстой. \\ \hline
NFR-SC2 & AI боловсруулах хэсгийг цаашид тусдаа сервис эсвэл queue-д суурилсан бүтэц рүү шилжүүлэх боломжтойгоор зохион байгуулах ёстой. \\ \hline
NFR-SC3 & Хичээл, тэмдэглэл, хавсралт файлын хэмжээ өсөхөд өгөгдлийн бүтэц болон хадгалалтын шийдэл шаталсан өргөтгөл даах боломжтой байх ёстой. \\ \hline
\end{tabular}
\caption{Өргөтгөх чадварын функцийн бус шаардлагууд}
\end{table}

\subsection{Найдвартай байдал}

\begin{table}[H]
\centering
\begin{tabular}{|p{3cm}|p{11cm}|}
\hline
\textbf{ID} & \textbf{Шаардлага} \\ \hline
NFR-R1 & Материал боловсруулах үед алдаа гарсан тохиолдолд систем хэрэглэгчид ойлгомжтой төлөв харуулж, дахин оролдох боломж олгох ёстой. \\ \hline
NFR-R2 & Хэрэглэгчийн хичээл, тэмдэглэл, шалгалтын үр дүнг тогтмол нөөцлөх боломжтой байх ёстой. \\ \hline
NFR-R3 & Гадаад AI үйлчилгээ түр хугацаанд доголдсон үед үндсэн өгөгдөл алдагдахгүй, урсгалыг дараа нь сэргээх боломжтой байх ёстой. \\ \hline
\end{tabular}
\caption{Найдвартай байдлын функцийн бус шаардлагууд}
\end{table}

\subsection{Ашиглахад хялбар байдал}

\begin{table}[H]
\centering
\begin{tabular}{|p{3cm}|p{11cm}|}
\hline
\textbf{ID} & \textbf{Шаардлага} \\ \hline
NFR-U1 & Интерфэйс нь монгол хэл дээр ойлгомжтой нэршилтэй, desktop болон mobile төхөөрөмж дээр зохицон ажиллах responsive загвартай байх ёстой. \\ \hline
NFR-U2 & Хэрэглэгч файл оруулсны дараах боловсруулалтын үе шатыг алхам бүрээр нь ойлгож чадахуйц визуал төлөвтэй байх ёстой. \\ \hline
NFR-U3 & Ахиц, шалгалтын үр дүн, хичээлийн төлөв зэргийг нэг дэлгэц дээр төвөггүй ойлгохуйц байдлаар харуулах ёстой. \\ \hline
\end{tabular}
\caption{Ашиглахад хялбар байдлын функцийн бус шаардлагууд}
\end{table}

\subsection{Засвар үйлчилгээ}

\begin{table}[H]
\centering
\begin{tabular}{|p{3cm}|p{11cm}|}
\hline
\textbf{ID} & \textbf{Шаардлага} \\ \hline
NFR-M1 & Систем нь frontend, backend, database, external AI integration гэсэн үүргээрээ тусгаарлагдсан бүтэцтэй байх ёстой. \\ \hline
NFR-M2 & Кодын бүтэц нь модульчлагдсан, шинэ боломж нэмэхэд бага нөлөөлөлтэй байх ёстой. \\ \hline
NFR-M3 & API болон өгөгдлийн бүтэц тодорхой баримтжуулалттай, багийн хамтын ажиллагаанд ойлгомжтой байх ёстой. \\ \hline
\end{tabular}
\caption{Засвар үйлчилгээний функцийн бус шаардлагууд}
\end{table}
